\paragraph{Theorem D.2.34 (Proved: every positive weight on the original Engel profile).}
On the $IV^*,III,I_1$ rational elliptic surface, let $G$
generate the Mordell--Weil group, retain
$M_S=\mathcal O(G-O)$, and use translation $6dG$ with
linearization zero $dI_1$, for any integer $d>0$.
Use the smooth $d$-period hexagonal Mumford cap at $I_1$.
At either additive site one may retain the original cap
or use any free finite logarithmic cap described below.
This gives an unbounded initial-weight family outside
the finite sphere-candidate catalogue.

For logarithmic caps put
\[
 m_0=3k,\quad m_1=4\ell,\qquad
 v_i=r_i\delta+n_i u_i,\qquad k,\ell>0.
\]
The exact freeness conditions are
\[
\begin{aligned}
 t_0&=\gcd(m_0,n_0,3r_0),&
 &\gcd(m_0,n_0,r_0)=1,\quad 3\mid m_0/t_0,\\
 t_1&=\gcd(m_1,n_1,2r_1),&
 &\gcd(m_1,n_1,r_1)=1,\quad 4\mid m_1/t_1 .
\end{aligned}                                             \tag{D.WE1}
\]
There is no coprimality restriction on $d$.
All noninvariant parts of affine translations are
included through the marked translation conjugacy.
At an original site set its effective Seifert data to
$\overline m_i=1,\overline n_i=0$; otherwise set them
to $m_i,n_i$. Define
\[
 P=\overline m_1\overline n_0+
   \overline m_0\overline n_1,\qquad
 \mathbb Z^r\oplus A_1=
 \operatorname{coker}
 \begin{pmatrix}
  -\overline n_0&\overline m_0\\
  -\overline n_1&-\overline m_1
 \end{pmatrix}.                                           \tag{D.WE2}
\]
The full fundamental group is this abelian group, and
$r=1$ precisely when $P=0$.

The exact degree-two matrix has target $(s,t,t_0',t_1')$.
For two logarithmic caps it is $B=(B_0\ B_1)$, where
\[
 B_0=
 \begin{pmatrix}
 -2dn_0&-6dr_0&0\\
 -k&0&r_0\\
 0&m_0&n_0\\
 0&0&0
 \end{pmatrix},
 \qquad
 B_1=
 \begin{pmatrix}
 -3dn_1&-6dr_1&0\\
 2\ell&0&-r_1\\
 0&0&0\\
 0&m_1&n_1
 \end{pmatrix}.                                           \tag{D.WE3}
\]
For an original $IV^*$ or $III$ cap replace its three
columns, respectively, by
\[
 O_0=
 \begin{pmatrix}2d&0&0\\0&1&0\\0&0&1\\0&0&0\end{pmatrix},
 \qquad
 O_1=
 \begin{pmatrix}3d&0&0\\0&-1&0\\0&0&0\\0&0&1\end{pmatrix}.
                                                               \tag{D.WE4}
\]
Let $T$ be the torsion of $\operatorname{coker}B$,
let $L$ count original additive caps, and put
$\beta=5$ if $IV^*$ is original, and $\beta=0$ otherwise.
Set
\[
 (b_2,b_3)=
 \begin{cases}
 (d-1+2r,\;2r),&L=0,\\
 (d-1+\beta+L+r,\;2\beta+2L),&L>0.
 \end{cases}                                               \tag{D.WE5}
\]
Then all the integral homology is
\[
\begin{array}{c|c}
j&H_j(X;\mathbb Z)\\\hline
0,6&\mathbb Z\\
1&\mathbb Z^r\oplus A_1\\
2&\mathbb Z^{b_2}\oplus T\\
3&\mathbb Z^{b_3}\oplus T\\
4&\mathbb Z^{b_2}\oplus A_1\\
5&\mathbb Z^r .
\end{array}                                                \tag{D.WE6}
\]
In particular $\chi(X)=2d$ for every choice.

There are useful closed forms. If both sites are
logarithmic and $\gcd(m_i,n_i)=1$, put
$g=\gcd(d,k,2\ell)$. Then
\[
 T=
 \begin{cases}
 \mathbb Z/g\oplus\mathbb Z/(d|P|/g),&P\ne0,\\
 \mathbb Z/g,&P=0.
 \end{cases}                                               \tag{D.WE7}
\]
The fundamental group is $\mathbb Z/|P|$ for $P\ne0$
and $\mathbb Z$ for $P=0$. With at least one original
cap the torsion groups, for every free remaining twist,
are
\[
\begin{array}{c|c}
\text{original sites}&T\\\hline
IV^*\text{ only}&
 \mathbb Z/(d\gcd(2,n_1))\oplus\mathbb Z/\gcd(m_1,n_1)\\
III\text{ only}&
 \mathbb Z/(d\gcd(3,n_0))\oplus\mathbb Z/\gcd(m_0,n_0)\\
\text{both}&\mathbb Z/d .
\end{array}                                                \tag{D.WE8}
\]
These are direct-sum descriptions, not necessarily
ordered invariant factors.

\begin{proof}
\emph{The weighted analytic construction.}
Engel's linearization identity for $M_S$ and $6G$
has base degree one
\cite[Lemma 2.1 and Remark 8.5]{manEngel26}.
Iterate that contracting generator $d$ times.
Translations preserve the base, so its linearization
is the product of $d$ copies and has divisor $dI_1$.
Equivalently, replace the long gamma period by $d$
times itself. On the punctured family this is the
ordinary torus cover defined by that sublattice.
The good-reduction models are preserved by it:
the integral matrices below have orders three and
four for every $d$. After the respective root base
changes one can therefore make the same finite
affine quotients as in D.2.5 and D.2.29.

At the cusp retain the same unimodular triangular
fan, now modulo the sublattice generated by one
short-direction period and $d$ times the other.
Its quotient tube is smooth and proper by the same
finite-period toric construction as D.1.2. Its
normalization has $d$ copies of $dP_6$ and
\[
 H_*(N;\mathbb Z)
 =(\mathbb Z,\mathbb Z^2,\mathbb Z^{d+3},
                         \mathbb Z^2,\mathbb Z^d).         \tag{D.WE9}
\]
The normalization proof of D.1.2 uses
$\operatorname{diag}(1,t-1,t-1)$ over
$\mathbb Z[t]/(t^d-1)$ and is valid for every positive
$d$, including $2$ and $3$.
This constructs the weighted family directly.
It does not assume that the filled threefold is a
free cover of a weight-one finite quotient when $d$
shares a factor with its multiplicities.

\emph{Integral exterior and finite-cap data.}
Use the basis $(\delta,a,b,\gamma)$ of D.2.5, with
the new gamma period. The actual monodromies are
\[
\begin{aligned}
 A_0&=\begin{pmatrix}
 1&-1&0&2d\\0&-1&-1&6d\\0&1&0&0\\0&0&0&1
 \end{pmatrix},\\
 A_1&=\begin{pmatrix}
 1&0&1&-3d\\0&0&1&0\\0&-1&0&6d\\0&0&0&1
 \end{pmatrix},&
 A_\infty&=\begin{pmatrix}
 1&0&0&d\\0&1&1&0\\0&0&1&0\\0&0&0&1
 \end{pmatrix}.
\end{aligned}                                             \tag{D.WE10}
\]
Their product is the identity. The two fixed lattices
are generated by $\delta$ and
\[
 u_0=\gamma+2d(a+b),\qquad
 u_1=\gamma+3d(a+b).
\]
The finite coinvariants still have
$b=-a,\delta=-3a$ and $b=-a,\delta=-2a$,
respectively, and $u_i$ maps to the primitive gamma
class. Thus their fixed-quotient indices and root
characters are exactly those of D.2.29, independently
of $d$. Its proof gives (D.WE1), including every
nonprimitive free numerator, the torsion orders
$t_i$ of the caps and their normal lines, and the
three actual Thom generators.

Let $V$ be the regular bundle over a graph joining
the three peripheral circles. The common fibre
coinvariant in degree two is free of rank one:
its primitive functional is
\[
 \xi_d=a^*\wedge b^*+6d\gamma^*\wedge\delta^* .
\]
The integral relations are
$\delta a=\delta b=a\gamma=b\gamma=0$ and
$\delta\gamma=-6d\,ab$, so $s=[ab]$ is primitive.
For the graph part the five columns
\[
 (-\delta,-\delta),\quad(b,-a),\quad
 (\delta,0),\quad(u_0,0),\quad(0,u_1)                     \tag{D.WE11}
\]
are an integral basis of
$\ker[(A_0-I)\ (A_1-I)]$.
They lie in the kernel and have a constant unit
$5$-row minor, using rows $0,2,3,4,7$ in zero-based
indexing. Their rational rank is five, so that unit
minor proves saturation for every $d$.
The common exterior-square relations likewise contain
a unit minor in the displayed primitive kernel of
$\xi_d$; the certificate gives its column indices
as a polynomial identity in $d$.

The first two columns of (D.WE11) are the actual
cusp normal tori in the angular directions
$\delta,a$. Their fibre coefficients are zero in
the rigidified cusp marking. Quotienting by these
two columns leaves the free target
$(s,t,t_0',t_1')$ used above.
The norms and contractions are
\[
\begin{aligned}
 N_0a&=-k\delta,&N_0\gamma&=m_0u_0,&
 \xi_d(a,u_0)&=2d,\\
 N_1a&=-2\ell\delta,&N_1\gamma&=m_1u_1,&
 \xi_d(a,u_1)&=3d,\\
 &&\xi_d(\gamma,\delta)&=6d .
\end{aligned}
\]
A covering-torus loop $\lambda$ has normal column
with graph part $N_i\lambda$ and fibre part
$-\xi_d(\lambda,v_i)$; the affine meridian has
graph part $v_i$ and fibre part zero. These are the
full Thom generators, including the torsion
character case of D.2.29. Since the delta sweep
at the second site is the negative of the first
modulo (D.WE11), they give exactly (D.WE3).

For an original additive cap, $Q=6dG$ is narrow and
the component Chern vector of $M_S$ is the same
primitive $e_G-e_O$ as at $d=1$. D.2.32's actual
principal two-torus model applies with the marked
logarithm multiplied by $d$. The regular elliptic
area covectors, in the order
$(\delta a,\delta b,\delta\gamma,ab,a\gamma,b\gamma)$,
are
\[
 (0,0,0,1,-2d,2d),\qquad
 (0,0,0,1,-3d,3d).
\]
Their full kernels have $\xi_d$ images $2d\mathbb Z$
and $3d\mathbb Z$, and both complete fixed normal
lattices have zero fibre correction. Thus (D.WE4)
is the exact integral replacement. The local free
cokernels have ranks $7-2=5$ and $2-2=0$, and each
original loses one first-homology rank, giving
one additional free bridge. These statements retain
the primitive Chern class rather than making a
zero-degree approximation.

\emph{The cusp boundary torsion and the component meridians.}
Write $E=X\setminus\operatorname{int}N$ and $M=\partial N$.
For $d>1$, $H_1(M)$ has an actual $\mathbb Z/d$ summand:
the monodromy kills $a$ and $d\delta$, whereas the
cusp kills both primitive angular circles. It would
therefore be incorrect to reuse a free-boundary
Mayer--Vietoris splitting from weight one.
Instead use the exact sequence of the pair $(X,E)$.

The exterior-square cusp image is primitive of rank
two for every $d$, so $H_2(M)$ is free of rank six.
Its two angular normal tori bound by their toric
section-disc sweeps. They form a primitive sublattice
because their Wang projections are the primitive
fixed basis $\delta,a$.
Furthermore
$H_3(N,M)=H^3(N)=\mathbb Z^2$ by (D.WE9).
Exactness bounds the kernel of
$H_2(M)\to H_2(N)$ by rank two. Hence those two
angular tori are its entire integral kernel.

Attaching the two finite or original additive caps
to $V$ first shows
\[
 H_2(E)/j_*\ker(H_2(M)\to H_2(N))
 \simeq\operatorname{coker}B\oplus\mathbb Z^{\beta+L}.
                                                               \tag{D.WE12}
\]
Indeed every finite boundary-to-cap degree-two map
is onto, including when the cap and normal line
have torsion; its kernel is the full Thom image.
At an original cap split its free local cokernel.
The degree-one sequence then has a free kernel of
rank $L$, since the first Betti number of $E$ is
one and each original lowers the target rank by one.
The angular cusp classes come from $H_2(V)$, so
they have zero image in this new bridge. Thus
quotienting by them preserves its free splitting.
This proves (D.WE12) without dropping boundary
torsion or assuming a torsion extension splits.

Now $H_2(N,M)=H^4(N)=\mathbb Z^d$, with its basis
given by small positive transverse discs to the
$d$ components. Their boundary classes in $H_1(E)$
are \emph{identical}, not merely proportional.
In the smooth height-one fan a component ray
$(v,1)$ has its primitive meridian $(v,1)$ in
angular-plus-base coordinates. Two such meridians
differ by integral multiples of $a,\delta$.
Both angular classes already vanish in $H_1(E)$:
the unchanged vertical columns of $A_0-I,A_1-I$
generate the entire vertical lattice primitively.
Thus every disc has boundary the same class
$x_0+x_1$, with the common boundary orientation;
at an original site its $x_i$ is zero.

The abelian presentation of $H_1(E)$ has generators
$\gamma,x_0,x_1$ and relations
$\overline m_i x_i=\overline n_i\gamma$.
It has free rank one. Its element $x_0+x_1$ has
infinite order precisely when $P\ne0$.
Therefore the kernel of the actual map
\[
 \mathbb Z^d\longrightarrow H_1(E),\qquad
 (z_1,\ldots,z_d)\longmapsto
                 (z_1+\cdots+z_d)(x_0+x_1)
                                                               \tag{D.WE13}
\]
is free of rank $d-1+r$. This remains true when
$H_1(E)$ or the indicated meridian has torsion.
The pair sequence and (D.WE12) now give a short
exact sequence with this free quotient, hence a
split sequence of abelian groups.

The matrix (D.WE3) has rational corank $r$.
For example a left-kernel vector has its other
three coordinates fixed by its $s$ coordinate;
the remaining equation is $dP$ times that coordinate.
If an original is present, (D.WE4) kills that
coordinate and the entire rational cokernel, so
the corank is zero. These facts and (D.WE13) give
the claimed $b_2$ and torsion $T$.

Van Kampen, separately, proves the group assertion.
The cusp kills the primitive $a,\delta$ loops;
finite monodromy then kills $b$. Only the central
$\gamma$ and the two meridians remain, with
$x_0x_1=1$ and the two power relations in (D.WE2).
Thus the full group is abelian with that exact
integer presentation. Finally the cusp has Euler
number $2d$ by (D.WE9), and every other cap has
Euler number zero. Integral duality and universal
coefficients give (D.WE5)--(D.WE6).

\emph{Closed Smith forms.}
When both numerators are primitive, the same
Bezout reductions as D.2.5 eliminate two unit
columns of (D.WE3) and leave
\[
 \begin{pmatrix}
 -2dn_0&-3dn_1\\-k&2\ell
 \end{pmatrix},\qquad \det=-dP .
\]
Its entries have gcd $\gcd(d,k,2\ell)$.
To check this prime by prime, a prime dividing
$k$ cannot divide $n_0$, one dividing $\ell$
cannot divide $n_1$, $3\nmid n_0$, and $2\nmid n_1$.
These observations give precisely the stated gcd
at $2,3$ and at every other prime. Its Smith
form gives (D.WE7), including its rank-one case.
For an original $IV^*$, eliminate $t,t_0'$ using
(D.WE4). The remaining relations are
$2ds=3dn_1s=0$ and
$m_1t_1'=n_1t_1'=0$, because $6dr_1s$ is already
zero. This is the first row of (D.WE8).
The original $III$ case is identical with
$3d$ and $2dn_0$, and with both originals the
remaining fibre relation is $\gcd(2d,3d)s=ds=0$.
\end{proof}

For two logarithmic caps, a simply connected member
has $P=\pm1$, which forces both local numerators
primitive and $\gcd(3k,4\ell)=1$. Thus $g=1$ and
\[
 H_2=\mathbb Z^{d-1}\oplus\mathbb Z/d,\qquad
 H_3=\mathbb Z/d,\qquad H_4=\mathbb Z^{d-1}.                \tag{D.WE14}
\]
This includes even and three-divisible $d$, where the
free-cover reduction in D.1 is unavailable. At minimal
orders and $d$ prime to six it recovers D.1's groups.
The exact smooth classification beyond the cases
already settled there is not inferred from (D.WE14).
With both original caps, (D.WE6) instead gives
\[
 \pi_1=\mathbb Z,\qquad
 H_2=\mathbb Z^{d+7}\oplus\mathbb Z/d,\qquad
 H_3=\mathbb Z^{14}\oplus\mathbb Z/d,\qquad
 H_4=\mathbb Z^{d+7},\quad H_5=\mathbb Z.
\]
At $d=1$ this is exactly K13; for every $d$ it is the
actual iterated-translation family.

The executable \path{code/manifolds_weighted_engel.py}
accepts arbitrary integer input, checks (D.WE1), and
returns the full group and homology presentation.
It verifies polynomial identities for the monodromies,
primitive exterior-square quotient and graph basis,
and compares the unreduced graph and full area kernels
with (D.WE3)--(D.WE4). Its retained output includes
nonprimitive free twists, nontrivial cap torsion,
every original/log choice and zero-Euler cases.
